%%%%%%%%%%%%%%%%%%%%%%%%%%%%%%%%%%%%%%%%%%%%%%%%%%%%%%%%%%%%%%%%%%%%%%%%%%%%%%
%
% 03_state_space.tex
%
%%%%%%%%%%%%%%%%%%%%%%%%%%%%%%%%%%%%%%%%%%%%%%%%%%%%%%%%%%%%%%%%%%%%%%%%%%%%%%

\begin{mosbox}{2. The Cognitive State Space}

\BodyFont

The Mathematical Operating System is modeled by the triple

\[
\boxed{
M=(C,\mathcal{F},s)
}
\]

where

\begin{center}

\begin{tabular}{p{2.5cm}p{13.5cm}}

$\mathbf{C}$ &
A finite stratified simplicial complex encoding higher-order conceptual
relationships.

\\[3mm]

$\mathbf{\mathcal F}$ &
A cellular sheaf assigning local mathematical information to every simplex
together with compatible restriction maps.

\\[3mm]

$\mathbf{s}$ &
A distinguished global section representing a coherent cognitive state.

\end{tabular}

\end{center}

Rather than storing isolated pieces of information,
MOS represents cognition as a structured mathematical object in which
local information, topology and global compatibility coexist within a
single framework.

\end{mosbox}

%%%%%%%%%%%%%%%%%%%%%%%%%%%%%%%%%%%%%%%%%%%%%%%%%%%%%%%%%%%%%%%%%%%%%%%%%%%%%%
%% VISUAL DECOMPOSITION
%%%%%%%%%%%%%%%%%%%%%%%%%%%%%%%%%%%%%%%%%%%%%%%%%%%%%%%%%%%%%%%%%%%%%%%%%%%%%%

\vspace{5mm}

\begin{center}
\begin{tikzpicture}[scale=1.0, every node/.style={font=\small}]
% Vertices of a regular tetrahedron projection
\coordinate (v1) at (-3,0);
\coordinate (v2) at (0,2);
\coordinate (v3) at (3,0);
\coordinate (v4) at (0,-2);
% Simplices (filled triangles) highlighting 2-simplices
\draw[simplex, thick] (v1)--(v2)--(v3)--cycle; % base triangle
\draw[simplex, thick] (v1)--(v4)--(v3); % opposite triangle
\draw[simplex, thick] (v2)--(v4); % edge connecting top and bottom
% Vertices
\foreach \p in {v1,v2,v3,v4}{
  \fill[MOSRed] (\p) circle (4pt);
}
% Labels for components
\node at (-3,1.2) {\Large Base Complex};
\node at (0,3) {\Large Cellular Sheaf};
\node at (3,1.2) {\Large Global Section};
% Arrows indicating flow
\draw[arrow] (-1,1.5) -- (1,1.5);
\draw[arrow] (1,-1.5) -- (3,-1.5);
\end{tikzpicture}
\end{center}

\begin{center}
\begin{tikzpicture}[scale=1.05]

%%%%%%%%%%%%%%%%%%%%%%%%%%%%%%%%%%%%%%%%%%%%%%%%%%%%%
%% COMPLEX
%%%%%%%%%%%%%%%%%%%%%%%%%%%%%%%%%%%%%%%%%%%%%%%%%%%%%

\node at (-9,6)
{\Large\bfseries Base Complex};

\coordinate (A) at (-10,2);

\coordinate (B) at (-6,2);

\coordinate (C) at (-8,5);

\coordinate (D) at (-8,3.5);

\draw[simplex] (A)--(B)--(C)--cycle;

\draw[simplex] (A)--(D);

\draw[simplex] (B)--(D);

\draw[simplex] (C)--(D);

\foreach \p in {A,B,C,D}
{

\fill[MOSOrange]
(\p)
circle(4pt);

}

%%%%%%%%%%%%%%%%%%%%%%%%%%%%%%%%%%%%%%%%%%%%%%%%%%%%%
%% ARROW
%%%%%%%%%%%%%%%%%%%%%%%%%%%%%%%%%%%%%%%%%%%%%%%%%%%%%

\draw[arrow]
(-4,3.5)
--
(-1,3.5);

%%%%%%%%%%%%%%%%%%%%%%%%%%%%%%%%%%%%%%%%%%%%%%%%%%%%%
%% SHEAF
%%%%%%%%%%%%%%%%%%%%%%%%%%%%%%%%%%%%%%%%%%%%%%%%%%%%%

\node at (2,6)
{\Large\bfseries Cellular Sheaf};

\node[concept] (v1) at (0,5)
{$F(v_1)$};

\node[concept] (v2) at (4,5)
{$F(v_2)$};

\node[concept] (e1) at (2,2.5)
{$F(e)$};

\draw[arrow]
(v1)
--
(e1);

\draw[arrow]
(v2)
--
(e1);

%%%%%%%%%%%%%%%%%%%%%%%%%%%%%%%%%%%%%%%%%%%%%%%%%%%%%
%% ARROW
%%%%%%%%%%%%%%%%%%%%%%%%%%%%%%%%%%%%%%%%%%%%%%%%%%%%%

\draw[arrow]
(5.5,3.5)
--
(8.5,3.5);

%%%%%%%%%%%%%%%%%%%%%%%%%%%%%%%%%%%%%%%%%%%%%%%%%%%%%
%% GLOBAL SECTION
%%%%%%%%%%%%%%%%%%%%%%%%%%%%%%%%%%%%%%%%%%%%%%%%%%%%%

\node at (12,6)
{\Large\bfseries Global Section};

\node[concept] (g1) at (10,5)
{$x_1$};

\node[concept] (g2) at (14,5)
{$x_2$};

\node[concept] (g3) at (12,2.5)
{$x_e$};

\draw[arrow]
(g1)
--
(g3);

\draw[arrow]
(g2)
--
(g3);

\node at (12,.8)
{\Large Compatible Assignments};

\end{tikzpicture}

\end{center}

%%%%%%%%%%%%%%%%%%%%%%%%%%%%%%%%%%%%%%%%%%%%%%%%%%%%%%%%%%%%%%%%%%%%%%%%%%%%%%
%% INTERPRETATION
%%%%%%%%%%%%%%%%%%%%%%%%%%%%%%%%%%%%%%%%%%%%%%%%%%%%%%%%%%%%%%%%%%%%%%%%%%%%%%

\begin{highlightbox}

\BodyFont

The three components of the state space play distinct mathematical roles.

\begin{itemize}

\item The simplicial complex specifies the combinatorial structure of
memory.

\item The cellular sheaf assigns local mathematical data to every cell.

\item A global section corresponds to a configuration whose local
assignments remain mutually compatible under every restriction map.

\end{itemize}

Consequently,

\[
\boxed{
\text{Geometry}
\longrightarrow
\text{Organization}
}
\]

\[
\boxed{
\text{Sheaf}
\longrightarrow
\text{Local Information}
}
\]

\[
\boxed{
\text{Global Section}
\longrightarrow
\text{Coherent Memory}
}
\]

\end{highlightbox}

%%%%%%%%%%%%%%%%%%%%%%%%%%%%%%%%%%%%%%%%%%%%%%%%%%%%%%%%%%%%%%%%%%%%%%%%%%%%%%
%% SMALL REMARK
%%%%%%%%%%%%%%%%%%%%%%%%%%%%%%%%%%%%%%%%%%%%%%%%%%%%%%%%%%%%%%%%%%%%%%%%%%%%%%

\begin{remarkbox}

\BodyFont

This decomposition is the mathematical foundation upon which the remaining
components of MOS are built. Subsequent sections introduce coherence
measures, semantic geometry, diagnostics and dynamics using this common
state-space representation.

\end{remarkbox}